%! Carrying Out a Test for the Difference of Two Population Means — Unit 7, Lesson 7.9 — Exit Ticket (Answer Key)
\documentclass[10pt]{article}
\usepackage{apstats-article}
\usepackage{apstats-key}   % pulls in -boxes; do NOT also load -boxes
\newcommand{\TallMath}[1]{$\displaystyle #1\rule[-1.4em]{0pt}{3.2em}$}

\begin{document}

\pageheader{Unit 7, Lesson 7.9}{Exit Ticket --- Answer Key}
\namedateperiod

\vspace{0.15in}

Sleep researchers compared nightly sleep duration for college athletes and non-athletes.
Summary statistics: Athletes ($n_1 = 30$, $\bar x_1 = 7.4$ h, $s_1 = 0.9$ h);
Non-athletes ($n_2 = 28$, $\bar x_2 = 6.8$ h, $s_2 = 1.1$ h).
Technology gives $t \approx 2.30$, $df \approx 53$, $p \approx 0.025$ (one-sided, $H_a: \mu_1 > \mu_2$).
All conditions have been verified.

\begin{enumerate}
  \item Interpret the $p$-value in context.
        \ansline{Assuming the true mean sleep duration is the same for athletes and non-athletes, there is about a 2.5\% probability of observing a difference in sample means as large as 0.6 hours (athletes higher) just by chance.}
        \ansline{}
        \ansline{}

  \item At $\alpha = 0.05$, write a complete conclusion in context.
        \ansline{Because $p \approx 0.025 < \alpha = 0.05$, we reject $H_0$. We have convincing evidence that the true mean nightly sleep duration for college athletes is greater than for non-athletes.}
        \ansline{}
        \ansline{}
\end{enumerate}

\end{document}
